% ARCHIVED FROM v82 EXACT SOURCE EXCERPT: 02_fixed-h-mcfg/japanese/main_JP.tex
% Source reference: archive/mcfg-v82-pretrim-full-20261018 (commit bcc0caf0eade66f552127171c0de24d2b92950b0)
% Restoration: replace the corresponding trimmed location with this excerpt (without these header lines).

\begin{proposition}[固定窓では括弧付きCFG族のファンアウト2代入可能性を回復できない]
\label{prop:fixed-window-doubling-fanout-two}
$\Gamma_{\mathrm{br}}=\{\mathtt a,\lfloor,\rfloor,\lceil,\rceil\}$ とし，
関連CFG原稿~\cite{kuriyamaCFG} の言語族 $L_n$ を
\[
 E_0\to\mathtt a,\qquad
 E_i\to\lfloor E_{i-1}E_{i-1}\rfloor
       \mid\lceil E_{i-1}\rceil\quad(1\le i\le n)
\]
の開始非終端記号 $E_n$ から生成する．
$h_{\mathrm U}$ は $\lceil,\rceil$ を2元半束の吸収元に，
それ以外の文字を単位元に写す準同型とする．
任意の $k,\ell\ge0$ および
$n\ge\max\{3,k+\ell+1\}$ について
\[
 L_n\notin\Cmcf{2}{h_{\Gamma_{\mathrm{br}}}^{k,\ell}},
 \qquad
 L_n\notin\Cmcf{2}{(h_{\Gamma_{\mathrm{br}}}^{k,\ell},h_{\mathrm U})}.
\]
したがって，$h_{\mathrm U}$ と組み合わせても，
\emph{固定}接頭辞--接尾辞窓を加えるだけでは
この族のファンアウト2での代入可能性は回復しない．
\end{proposition}
\begin{proof}
$W_0=\mathtt a$，
$W_i=\lfloor W_{i-1}W_{i-1}\rfloor$ とおく．
$L_n$ のうち天井括弧を含まない唯一の語は $W_n$ である．
完全二分木の語 $W_{n-1}$ を
$W_{n-1}=B\mathtt a R$，
$B=\lfloor^{n-1}$ と書けば
\[
 W_n=\lfloor B\mathtt a R\,B\mathtt a R\rfloor
\]
となる．$B,R$ はいずれも非空で，
$|B|=n-1\ge k+\ell$，
$|R|=3\cdot2^{n-1}-n-2\ge n-1\ge k+\ell$ である．

次の2タプルと2つの文文脈を考える：
\[
 \begin{aligned}
 \tuple x&=(B,\ R B\mathtt a R),&
 \tuple y&=(B\mathtt a R B,\ R),\\
 E&=\lfloor\blank_1\mathtt a\blank_2\rfloor,&
 F&=\blank_1\lfloor\mathtt a\blank_2\rfloor.
 \end{aligned}
\]
どちらの文脈でも2つの穴の間の文字列は非空であり，
定義~\ref{def:tuple-substitutable} に適合する．
第1成分の $B$ と $B\mathtt a R B$ は，
長さ $k$ の接頭辞と長さ $\ell$ の接尾辞が一致する．
第2成分の $R B\mathtt a R$ と $R$ も同様である．
各成分は $k+\ell$ 以上の長さを持ち，
天井括弧を含まない．
よって $\tuple x,\tuple y$ の成分別の型は，
主張にある両方の型付けで一致する．
しかし
\[
 E[\tuple x]=E[\tuple y]=F[\tuple x]=W_n,
 \qquad
 F[\tuple y]=B\mathtt a R B\lfloor\mathtt a R\rfloor
 \ne W_n.
\]
$F[\tuple y]$ の最初の $\mathtt a$ の前には
$\lfloor$ が $n-1$ 個しかないのに対し，
$W_n$ では $n$ 個あるためである．
$F[\tuple y]$ にも天井括弧は含まれないため，
$F[\tuple y]\notin L_n$ である．
従って両方の型付けについて
$(2,h)$-タプル代入可能性が破れる．
\end{proof}

\begin{remark}
本命題は元のCFG族が高ファンアウトの標的クラスに
\emph{属さない}ことを示すのであって，
ファンアウト2における特徴データ長の超多項式下界ではない．
特に問題~\ref{op:fanout-two} は未解決であり，
その解決には $\Cmcf{2}{h}$ に実際に属する
\emph{別の}言語族の構成が必要である．
\end{remark}



\subsection{希少な型と2成分の局所的可動性}
\label{subsec:rare-type-mobility}

\begin{lemma}[$h$ の型全体の削除]
\label{lem:mcfg-delete-type}
$f\ge1$、$L\in\Cmcf{f}{h}$、$m\in M$ に対し，
$L^{\setminus m}:=L\setminus h^{-1}(m)$ も
$\Cmcf{f}{h}$ に属する．
\end{lemma}
\begin{proof}
$h^{-1}(M\setminus\{m\})$ は正規言語であり，
$(f,h)$-タプル代入可能でもある．実際，成分別の型が
一致すれば，どの文文脈でも全体の型が一致する．
共通部分は代入可能性を保存し
（命題~\ref{prop:abelian-coset-filter}），
$f$-MCFL は正規言語との共通部分で閉じる
~\cite{SekiEtAl1991}．
\end{proof}

\begin{proposition}[希少な型は長い正例を強制する]
\label{prop:mcfg-rare-type-sample}
$L\in\Cmcf{f}{h}$、$m\in h(L)$ とし，
$\mu_m(L)=\min\{|w|:w\in L,\ h(w)=m\}$ とおく．
$\mathcal B^{\le r}_{f,h}$（$r\ge1$）または
$\mathcal B_{f,h}$ の任意の特徴標本には
型 $m$ の語が含まれ，従って符号化総長は
$\mu_m(L)+1$ 以上である．
より一般に，$L$ と $L^{\setminus m}$ の両方に
特徴標本を持つ任意の集合駆動型再構成器についても
同じ必要性が成立する．
\end{proposition}
\begin{proof}
型 $m$ を避ける標本 $K\subseteq L$ は
$K\subseteq L^{\setminus m}$ を満たす．
補題~\ref{lem:mcfg-delete-type} と
命題~\ref{prop:reconstruction-soundness} を
標的 $L^{\setminus m}$ に適用すると，
$K$ からの再構成言語は真部分言語
$L^{\setminus m}$ に含まれる．
一般の集合駆動型再構成器については，
$L$ の特徴標本 $\mathcal F$ が型 $m$ を避けると仮定し，
$L^{\setminus m}$ の特徴標本を $\mathcal F'$ とする．
有限集合 $\mathcal F\cup\mathcal F'$ は両標的に含まれ，
双方の特徴標本を含むので，同じ再構成結果が
異なる二言語に等しくなり矛盾する．
\end{proof}

\begin{lemma}[ファンアウト2の冪移送]
\label{lem:mcfg-power-transfer}
$f\ge2$、$L\in\Cmcf{f}{h}$、$u\in\Sigma^+$ とする．
$g=h(u)$ とし，ある $i_0,\pi\ge1$ が
全ての $i\ge i_0$ について
$g^{i+\pi}=g^i$ を満たすとする．
$s,r\ge i_0$、$d\in\pi\mathbb Z_{>0}$ とする．
$L$ のある語が $u^{s+d+r+2}$ を因子に持つなら，
任意の $X,Z\in\Sigma^*$、$Y\in\Sigma^+$ に対し
\[
 Xu^sYu^{d+r}Z\in L
 \quad\Longleftrightarrow\quad
 Xu^{s+d}Yu^rZ\in L
\]
が成り立つ．
\end{lemma}
\begin{proof}
受理語を $Pu^{s+d+r+2}Q$ と書く．
$\tuple x=(u^s,u^{d+r})$、
$\tuple y=(u^{s+d},u^r)$ とすると，
両者の成分別 $h$-型は一致する．
文脈 $E=Pu\blank_1u\blank_2Q$ は内部区切りが非空で，
$E[\tuple x]=E[\tuple y]=Pu^{s+d+r+2}Q\in L$ である．
よって2成分代入可能性から
両タプルの分布は一致する．
任意の適切な文脈 $X\blank_1Y\blank_2Z$
に適用すれば主張が従う．
\end{proof}

\begin{lemma}[有界な局所シフト対]
\label{lem:mcfg-local-shift}
$\Sigma$ は空でないとする．
$c=|M|^2|\Sigma|$ とし，$R_c(5)$ を
5頂点完全グラフに対する $c$ 色Ramsey数とする．
$|w|\ge R_c(5)$ を満たす任意の語は
$w=P\,x\,p\,q\,p\,y\,Q$ と分解でき，
$x,y\in\Sigma^+$、$p\in\Sigma$、
$q\in\Sigma^*$ および
\[
 h(x)=h(xpq),\qquad h(qpy)=h(y)
\]
を満たす．従って $w\in L\in\Cmcf{2}{h}$ ならば
$(x,qpy)$ と $(xpq,y)$ の分布は一致する．
\end{lemma}
\begin{proof}
最初の $R_c(5)$ 個のカット位置から $i<j$ を選び，
$(h(w[i:j]),h(w[i+1:j]),w[i])$ で辺を彩色する．
単色5点集合 $i_1<\cdots<i_5$ を取り，
$x=w[i_1:i_2]$、
$p=w[i_2]=w[i_3]$、
$q=w[i_2+1:i_3]$、
$y=w[i_3+1:i_5]$ とする．
$P,Q$ は残りの接頭辞・接尾辞とする．
辺 $(i_1,i_2),(i_1,i_3)$ の第一座標が
$h(x)=h(xpq)$ を与え，
$(i_2,i_5),(i_3,i_5)$ の第二座標が
$h(qpy)=h(y)$ を与える．
内部区切り $p$ が非空である文脈
$P\blank_1p\blank_2Q$ は両タプルを受理するため，
2成分代入可能性によって分布が一致する．
\end{proof}

\begin{theorem}[任意の固定有限型付けは括弧付き倍化族を救済できない]
\label{thm:mcfg-all-typing-doubling-obstruction}
任意の有限モノイド準同型
$h:\Gamma_{\mathrm{br}}^*\to M$ について，
命題~\ref{prop:fixed-window-doubling-fanout-two} の $L_n$ は
\[
 n\ge2|M|+3\quad\Longrightarrow\quad
 L_n\notin\Cmcf{2}{h}
\]
を満たす．$f\ge2$ に対しても
$L_n\notin\Cmcf{f}{h}$ となる．
\end{theorem}
\begin{proof}
$g=h(\lfloor)$ の冪の指数を $i_0\ge1$、
周期を $\pi\ge1$ とし，
$i_0+\pi\le|M|+1$、$i_0\le|M|$ とする．
$n\ge2i_0+\pi+2$ の場合を考え，
$s=i_0$、$d=\pi$、
$r=n-s-d-2\ge i_0$ とおく．
$W_n=\lfloor W_{n-1}W_{n-1}\rfloor$ は
$L_n$ に属する唯一の天井括弧なしの語で，
$\lfloor^n$ から始まる．
ゆえに $u=\lfloor$ として
補題~\ref{lem:mcfg-power-transfer} の前提が成立する．

同じ $W_n$ の先頭ランから，
最初の $\lfloor$ を除いた位置に
第一ブロック $\lfloor^s$ を取る．
第二ブロック $\lfloor^{d+r}$ は
二番目の $W_{n-1}$ の先頭ラン内に取る．
両者の間の文字列 $Y$ は最初の $\mathtt a$ を
含むため非空である．従って適切な $X,Z$ により
$W_n=X\lfloor^sY\lfloor^{d+r}Z$ と書ける．
冪移送から
$X\lfloor^{s+d}Y\lfloor^rZ\in L_n$
が従うはずだが，これは天井括弧を含まず，
最初の $\mathtt a$ の前の $\lfloor$ が
$n+d$ 個ある．一方，唯一の天井括弧なしの語
$W_n$ では $n$ 個であり，矛盾する．
$2i_0+\pi+2\le2|M|+3$ より一様な閾値が得られる．
\end{proof}

\begin{remark}
定理~\ref{thm:mcfg-all-typing-doubling-obstruction} は，
命題~\ref{prop:fixed-window-doubling-fanout-two} の
窓型付けに限った障害を任意の有限型付けへ強化する．
窓型付け固有の鋭い閾値は旧命題に残す．
これは言語の\emph{クラス非所属}の定理であり，
問題~\ref{op:fanout-two} が問う特徴データの
高ファンアウト下界ではない．
\end{remark}

窓型付けと任意の有限群を組み合わせた場合の拡張は，
\ref{app:window-group} で証明した．
可換群の場合のより小さい定数は
\ref{app:window-abelian} に残す．
固定窓のみの場合の二次上界と窓・有限群型付けの場合の
三次上界について，文法サイズの指数が最適かどうかは未解決である．
非可換群の定数 $M_{A,f}$ の $f$ への二重指数的依存を下げられるかも今後の課題である．
付録~\ref{app:local-group} は，正像が局所群となるすべての型付け，
すなわち2元半束を含まない最大のクラスに，この元文法厚みによる制御を拡張する．
対応する下界はファンアウト1でのみ知られている（系~\ref{cor:pseudovariety-dichotomy}）．
系~\ref{cor:mcfg-li-group-polydata} は，
$\mathbf{LI}\vee\mathbf{Grp}$ による見通しのよい十分条件を与える．
より広い局所群型付けと擬多様体レベルの二分法は
付録~\ref{app:local-group} で別に扱う．
高ファンアウトの下界問題は引き続き未解決である．

今後の方向性としては，
複数のファンアウトバイアスにまたがる適応的学習，
値域サイズの事前上界なしでの有用な型付けの学習，
型付きクラス内部の規則ランク1部分クラスの特徴づけ，
全多重文脈に基づく合同的 MCFG との関係の解明，
および PMCFG 型のコピーを許す再構成への拡張がある．

\section*{謝辞}
著者は有益な議論と励ましに対し，
金澤誠氏および吉仲亮氏に感謝する．

\section*{利益相反に関する申告}
著者は，本論文の研究に影響を及ぼしたように見え得る
金銭上の利益関係または個人的関係を認識していない．

\section*{研究資金}
本研究は公的機関，商業部門，非営利部門の
いずれの助成機関からも特定の研究助成を受けていない．

\section*{データの利用可能性}
実験的なデータセットは使用していない．
有限例の検算を再現するスクリプトと原稿ソースは
公開リポジトリ
\url{https://github.com/growupkuriyama-hub/Papers/tree/main/02_fixed-h-mcfg}
に置いた．一般定理の根拠は本論文の数学的証明であり，
有限計算の網羅性を一般証明とみなしてはいない．

\section*{原稿作成過程における生成AI・AI支援技術の利用申告}
著者は研究および原稿作成の過程で
ChatGPT（OpenAI）と Claude（Anthropic）を使用した．
用途には，数学的論証の探索，証明草稿の作成と批判的検討，
文献・整合性の確認，文章と LaTeX ソースの編集が含まれる．
AI が生成した提案や査読を，
独立した数学的検証とみなすものではない．
本論文の数学的主張，引用文献，最終的な内容については
著者が責任を負う．

\appendix
